\section{Capítulo~\ref{sec:muscles}. La salida a los músculos}

La estructura de la salida (unidades motoras, margen de seguridad de la sinapsis, activación por tamaño, par de músculos, lazo de estiramiento y generadores de ritmo) está medida directamente; la condición de estabilidad del lazo con retardo es un teorema; el modelo interno del cuerpo es una hipótesis bien fundada; los números de las figuras son cálculos con el modelo del libro.

{\footnotesize
\setlength{\tabcolsep}{3pt}\renewcommand{\arraystretch}{1.15}
\begin{longtable}{>{\raggedright}p{0.5cm}>{\raggedright}p{4.0cm}>{\raggedright}p{5.6cm}>{\raggedright}p{2.3cm}>{\raggedright\arraybackslash}p{1.7cm}}
\toprule
N.º & Afirmación & Experimento y qué se midió & Fuente & Estado \\
\midrule\endfirsthead
\toprule
N.º & Afirmación & Experimento y qué se midió & Fuente & Estado \\
\midrule\endhead
1 & Unidad motora: una neurona \nt{ACET} controla un grupo de fibras, de unos cientos a un par de miles; en los músculos de ajuste fino las unidades son más pequeñas & Músculos humanos, recuento de axones motores y de fibras musculares: en promedio unas $340$ fibras por neurona en un músculo interóseo de la mano, unas $410$ en el braquiorradial, unas $1930$ en el gastrocnemio medial. El capítulo no da una cifra exacta para las unidades más pequeñas. & \cite{Feinstein1955mu} & medido \\
2 & La sinapsis sobre el músculo libera varias veces más neurotransmisor del necesario para activar la fibra & Revisión de mediciones en sinapsis neuromusculares: el margen de seguridad (cociente entre el neurotransmisor liberado y el umbral) en mamíferos adultos suele ser de $3$--$5$; en humanos, el margen se apoya más en los pliegues de la membrana receptora que en la cantidad liberada. & \cite{WoodSlater2001} & medido \\
3 & La sacudida~\eqref{eq:twitch} sube en un tiempo $T_c$ del orden de $50\ms$ & Unidades motoras individuales de la mano humana con esfuerzo voluntario, fuerza promediada sobre las espigas de la unidad: tiempo de subida de $30$--$100\ms$, menos de $70\ms$ en el $80\,\%$ de las unidades. La función $(t/T_c)e^{1-t/T_c}$ es una aproximación tomada de un modelo del conjunto de unidades. & \cite{MilnerBrown1973rec, Fuglevand1993} & medido \\
4 & Suma de sacudidas~\eqref{eq:forcesum}: con $5$, $15$ y $40$ espigas/s, fuerza de $0.2$--$1.1F_1$, $1.8$--$2.2F_1$ y cerca de $5.4F_1$ (fig.~\ref{fig:tetanus}) & Cálculo con \texttt{models/\allowbreak muscle\_\allowbreak models.py}, $T_c=50\ms$: $0.21$--$1.10$, $1.76$--$2.19$ y $5.28$--$5.49\,F_1$ en los últimos $0.3\,\text{s}$. La suma lineal es una simplificación: el modelo no incluye la saturación del músculo real. & \texttt{models/\allowbreak muscle\_\allowbreak models.py} & modelo \\
5 & Las unidades se activan por tamaño; las más débiles son cien veces más débiles que las más fuertes & Raíces ventrales del gato: con una entrada refleja creciente, las primeras en descargar son las neuronas con espigas pequeñas en la raíz, es decir, las pequeñas. Músculo interóseo de la mano humana: fuerza de sacudida de las unidades de $0.1$ a $10$~g, que crece casi linealmente con la fuerza a la que se activa la unidad. & \cite{Henneman1957, MilnerBrown1973rec} & medido \\
6 & El paso de fuerza es proporcional a la fuerza, $\Delta F/F\approx q-1$~\eqref{eq:sizeprinciple}: punto flotante & Deducido en el capítulo a partir de la progresión geométrica de las fuerzas. Concuerda con los experimentos: con esfuerzos grandes, para el mismo incremento de fuerza se activan muchas menos unidades nuevas. & deducción en el capítulo; \cite{MilnerBrown1973rec} & teoría \\
7 & La dispersión de la fuerza crece en proporción a la propia fuerza & Esfuerzo isométrico voluntario de un músculo del pulgar: la desviación estándar de la fuerza crece linealmente con la fuerza media; con el mismo esfuerzo provocado por estimulación eléctrica, la dispersión es constante. Modelo del conjunto de unidades: el crecimiento lineal lo produce la activación de las unidades por orden de fuerza. & \cite{Jones2002sdn} & medido \\
8 & La suavidad de los movimientos es consecuencia del ruido dependiente de la señal & Modelo: la trayectoria que minimiza la dispersión del punto final con ese ruido reproduce las trayectorias medidas del ojo y del brazo. No hay un experimento directo que distinga esta causa de otras. & \cite{HarrisWolpert1998} & hipótesis \\
9 & La diferencia de fuerzas del par de músculos fija el momento; la suma, la rigidez~\eqref{eq:antagonists} & Teoría del control de la impedancia mediante cocontracción. Brazo humano: la rigidez de cada articulación, medida con pequeños empujones, es aproximadamente proporcional a su momento; distintas proporciones de cocontracción cambian la forma y la dirección de la elipse de rigidez de la mano. & \cite{Hogan1984, GomiOsu1998stiff} & medido \\
10 & El sensor de estiramiento excita su propia neurona \nt{ACET} y, a través de una intermediaria inhibidora, apaga al antagonista (fig.~\ref{fig:antagonists}) & Médula espinal del gato: una sola intermediaria, excitada por las fibras de los sensores de estiramiento, produce potenciales inhibidores monosinápticos en las neuronas \nt{ACET} del músculo antagonista, con un retardo sináptico de $0.28$--$0.42\ms$. El neurotransmisor de la intermediaria (glicina) no se determinó en este experimento. & \cite{Jankowska1972Ia} & medido \\
11 & Retardo del lazo: espinal, $25$--$30\ms$; a través de la corteza, entre una vez y media y tres veces más; a través del ojo, $100$--$200\ms$ & Brazo humano, empujón a la articulación: respuesta muscular corta a los $20$--$45\ms$, larga a los $45$--$105\ms$, voluntaria a los $120$--$180\ms$. Desplazamiento de la posición visible del dedo durante el movimiento: corrección a los $160\ms$ en promedio. & \cite{Pruszynski2008rapid, SaundersKnill2003vis} & medido \\
12 & $\dd x/\dd t=-Gx(t-d)$ es estable con $Gd<\pi/2$~\eqref{eq:delaystable}; en el límite la frecuencia es $1/(4d)$ & Teorema sobre las raíces de una ecuación con retardo. Comprobación con \texttt{models/\allowbreak muscle\_\allowbreak models.py}, $d=30\ms$: a los $3\,\text{s}$ la amplitud es $3\cdot10^{-6}$ con $G=40$, $0.13$ con $G=50$ y $38$ con $G=55$ por segundo (límite: $52$). & \cite{Hayes1950} & teoría \\
13 & Temblor fisiológico de $8$--$12$~Hz; el lazo de estiramiento es una de sus fuentes & Revisión de mediciones: las personas sanas tienen un temblor fino en torno a $10$~Hz; su origen es mixto (oscilaciones centrales, propiedades de las descargas de las unidades, resonancia mecánica y resonancia del lazo reflejo), y en distintas condiciones predomina uno u otro. & \cite{McAuleyMarsden2000} & hipótesis \\
14 & El cerebro predice el resultado de una orden con un modelo interno del cuerpo & Una persona sostiene una carga en pinza y mueve el brazo: con carga inercial, viscosa y elástica, la fuerza de prensión cambia al mismo tiempo que la carga, no con el retardo del lazo, es decir, la carga se predice. Una revisión reúne estos experimentos; no hay observación directa del modelo mismo en las neuronas. & \cite{FlanaganWing1997grip, WolpertGhahramani2000} & hipótesis \\
15 & El ritmo de caminar, respirar y masticar lo generan redes locales con una entrada constante & Revisión de experimentos: un sistema nervioso aislado (médula espinal, ganglios) produce un patrón motor rítmico sin entrada rítmica y sin realimentación de los músculos. & \cite{MarderBucher2001} & medido \\
16 & La inhibición mutua con fatiga~\eqref{eq:halfcentre} da un ritmo de período $0.84\,\text{s}\approx2\tau_a$ (fig.~\ref{fig:halfcentre}) & Teorema: una red con inhibición mutua y adaptación sin estado estable oscila necesariamente con entrada constante. Cálculo con \texttt{models/\allowbreak muscle\_\allowbreak models.py} ($I=1$, $\beta=2$, $b=2$, $\tau=20\ms$, $\tau_a=0.4\,\text{s}$): período de $0.84\,\text{s}$. No se ha demostrado que los generadores reales estén construidos así. & \cite{Matsuoka1985osc} & modelo \\
\bottomrule
\end{longtable}}
